\section{Ritos Scrum no calendário}
\label{sec:ritos}

Para cada sprint $k=1,\dots,K$, indexamos dias $d=1,\dots,D$ e um dia adicional formal de \emph{retrospectiva} sem alteração mecânica do estado no protótipo atual (apenas registro em log).

\begin{table}[htbp]
  \centering
  \caption{Mapeamento simples entre dia no sprint e modo de processamento.}
  \label{tab:ritos}
  \begin{tabular}{@{}llp{7.2cm}@{}}
    \toprule
    Dia $d$ & Rito / modo & Efeito no motor (v1) \\
    \midrule
    $1$ & \emph{Sprint Planning} & Puxa cartões de \texttt{backlog} para \texttt{analise} até $P_{\max}$ (e WIP); em seguida, se $P_{\max}>0$, continua a preencher \texttt{analise} a partir do \texttt{backlog} até ao WIP. \\
    $2,\dots,D-1$ & \emph{Daily Scrum} & Dia de trabalho: aplica capacidade, alocação e avanços. \\
    $D$ & \emph{Sprint Review} & Também dia de trabalho (mesma mecânica), com marcação explícita no log. \\
    $D{+}1$ & \emph{Retrospectiva} & Sem efeito mecânico; encerra sprint e reinicia $d\leftarrow 1$. \\
    \bottomrule
  \end{tabular}
\end{table}

A separação pedagógica entre ``marcador'' de Daily e o processamento de trabalho facilita narrativas de ensino: o calendário Scrum aparece como \textbf{ritmo} sobreposto ao \textbf{fluxo} Kanban~\cite{schwaber2002agile}, sem exigir modelagem de horas de reunião na primeira versão. Em termos de modelagem OR, os ritos funcionam como \textbf{modos} do mesmo relógio de simulação: mudam quais transições são habilitadas no início do dia (puxar \emph{vs.} apenas processar), preservando um único estado de sistema entre dias.

\paragraph{Eventos aleatórios de capacidade (Daily e Review).}
Nos dias de \emph{Daily Scrum} e \emph{Sprint Review}, o motor aplica uma probabilidade global configurável (\texttt{dailyRandomEventChance}): com essa probabilidade, sorteia-se \emph{no máximo} um evento de um catálogo ponderado (doença, férias, meio período, reunião de equipa, incidente de infraestrutura, ``bom clima''), cada qual com multiplicador ou anulação de capacidade efetiva para membros afetados. Os acontecimentos aparecem no \texttt{DayLog} e são listados na interface para consulta pedagógica.
